\documentclass[10pt, aspectratio=169]{beamer}
\input{../../_en_preambulo_beamer}
% Comandos y macros estadísticas compartidas para las presentaciones Beamer en inglés (Metropolis)

% Conjuntos numéricos básicos
\newcommand{\R}{\mathbb{R}}
\newcommand{\N}{\mathbb{N}}
\newcommand{\Q}{\mathbb{Q}}
\newcommand{\Z}{\mathbb{Z}}
\newcommand{\set}[1]{\left\{ #1 \right\}}
\newcommand{\abs}[1]{\left|#1\right|}
\newcommand{\norm}[1]{\left\| #1 \right\|}

% Comandos estadísticos y probabilísticos del curso
\newcommand{\Var}{\operatorname{Var}}
\newcommand{\cov}{\operatorname{Cov}}
\newcommand{\corr}{\rho}
\newcommand{\comb}[2]{\begin{pmatrix} #1 \\ #2 \end{pmatrix}}
\newcommand{\s}{\sigma}
\newcommand{\particion}{\mathfrak{P}}
\newcommand{\card}[1]{\#\left( #1 \right)}
\newcommand{\seq}[3]{\set{#1}_{#2}^{#3}}
\newcommand{\E}{\mathbb{E}}
\newcommand{\Prob}{\mathbb{P}}

% Entornos pedagógicos y cajas en inglés académico natural (soportando tanto nombres en inglés como en español)
\newenvironment{discussion}{%
  \begin{alertblock}{Discussion Question}%
}{%
  \end{alertblock}%
}
\newenvironment{discusion}{%
  \begin{alertblock}{Discussion Question}%
}{%
  \end{alertblock}%
}

\newenvironment{reflection}{%
  \begin{alertblock}{Classroom Reflection}%
}{%
  \end{alertblock}%
}
\newenvironment{reflexion}{%
  \begin{alertblock}{Classroom Reflection}%
}{%
  \end{alertblock}%
}

\newenvironment{paradox}{%
  \begin{alertblock}{Exploratory Paradox}%
}{%
  \end{alertblock}%
}
\newenvironment{paradoja}{%
  \begin{alertblock}{Exploratory Paradox}%
}{%
  \end{alertblock}%
}

\newenvironment{motivation}{%
  \begin{exampleblock}{Motivation: Data Science in Practice}%
}{%
  \end{exampleblock}%
}
\newenvironment{motivacion}{%
  \begin{exampleblock}{Motivation: Data Science in Practice}%
}{%
  \end{exampleblock}%
}

\newenvironment{quickchallenge}{%
  \begin{block}{Quick Team Challenge}%
}{%
  \end{block}%
}
\newenvironment{retoclase}{%
  \begin{block}{Quick Team Challenge}%
}{%
  \end{block}%
}


\title[Hypothesis Testing Foundations]{Section 07.01: Foundations of Hypothesis Testing}
\subtitle{$H_0$ vs. $H_1$, Type I and II Errors, and Statistical Power}
\author[J. Castillo Colmenares]{Juliho Castillo Colmenares}
\institute{Tecnológico de Monterrey}
\date{\vspace{-1.5cm}}

\begin{document}

% Slide 1: Title
\begin{frame}[plain]
  \titlepage
\end{frame}

% Slide 2: Roadmap
\begin{frame}{Roadmap --- Chapter 07: Hypothesis Testing}
  \footnotesize
  Where are we in the modular development of the chapter? \pause
  \begin{itemize}\itemsep=0.08em
    \item \textbf{07.01} Foundations: $H_0$ vs. $H_1$, Type I and II Errors, and Power \emph{(Today)} \pause
    \item \textbf{07.02} $Z$ and $t$ Tests for One- and Two-Sample Means \pause
    \item \textbf{07.03} Chi-Squared Goodness-of-Fit Tests \pause
    \item \textbf{07.04} Contingency Tables and Independence Tests
  \end{itemize}
  \vspace{-0.05cm}
  \begin{block}{Session Objective}
    Formalize the complete logical and probabilistic framework of hypothesis testing: how $H_0$ and $H_a$ are formulated, what Type I ($\alpha$) and Type II ($\beta$) errors are, what a test's power measures, and how to determine the sample size needed to reach a target power.
  \end{block}
\end{frame}

% Slide 3: Motivation
\begin{frame}{From Estimation to Decision}
  \footnotesize
  \begin{motivacion}
    In Chapter 06 we learned to \emph{estimate} population parameters via points and intervals. We now take the next logical step: using sample data to \emph{decide} between two opposing claims about those parameters --- the heart of applied statistical inference.
  \end{motivacion}

  \vspace{0.15cm}
  \pause
  \begin{itemize}\itemsep=0.1cm
    \item \textbf{Null hypothesis ($H_0$):} the standing claim, assumed true until evidence against it. \pause
    \item \textbf{Alternative hypothesis ($H_a$):} what we seek to establish; it negates $H_0$. \pause
    \item \textbf{Every decision is uncertain:} no test is infallible; we must quantify the risk of being wrong.
  \end{itemize}
\end{frame}

% Slide 4: Formal notation and test types
\begin{frame}{Formal Notation and Test Types}
  \footnotesize
  \begin{block}{Formulating Hypotheses}
    $$H_0: \mu = \mu_0 \qquad \text{versus} \qquad H_a: \mu \neq \mu_0 \ \text{(or } < \ \text{or } > \text{)}$$
  \end{block}
  \pause

  \vspace{0.1cm}
  \begin{alertblock}{Three Possible Structures}
    \begin{itemize}\itemsep=0.05cm
      \item \textbf{Left-tailed:} $H_a: \mu < \mu_0$.
      \item \textbf{Right-tailed:} $H_a: \mu > \mu_0$.
      \item \textbf{Two-tailed:} $H_a: \mu \neq \mu_0$.
    \end{itemize}
    The direction of $H_a$ always reflects the suspicion or claim being tested.
  \end{alertblock}
\end{frame}

% Slide 5: Type I and II Errors
\begin{frame}{The Two Types of Error in Hypothesis Testing}
  \footnotesize
  Every statistical decision based on a sample can go wrong in two distinct ways: \pause

  \vspace{0.1cm}
  \begin{block}{Type I Error ($\alpha$)}
    Rejecting $H_0$ when it is actually true.
    $$\alpha = P(\text{Reject } H_0 \mid H_0 \text{ true})$$
  \end{block}
  \pause

  \vspace{0.1cm}
  \begin{alertblock}{Type II Error ($\beta$)}
    Failing to reject $H_0$ when it is actually false.
    $$\beta = P(\text{Do not reject } H_0 \mid H_0 \text{ false})$$
  \end{alertblock}
\end{frame}

% Slide 6: Power and sample size formula
\begin{frame}{Power and Sample Size}
  \footnotesize
  \begin{block}{Power of the Test}
    $$\text{Power} = 1-\beta = P(\text{Reject } H_0 \mid H_0 \text{ false})$$
    It is the test's ability to detect a real effect when one truly exists.
  \end{block}
  \pause

  \vspace{0.1cm}
  \begin{alertblock}{Sample Size for a Target Power (One-Tailed $Z$-Test)}
    $$n = \left(\frac{(Z_\alpha+Z_\beta)\sigma}{\mu_a-\mu_0}\right)^2$$
    Higher required power or a smaller tolerated $\alpha$ increases the required $n$; a larger effect size ($\mu_a-\mu_0$) decreases the required $n$.
  \end{alertblock}
\end{frame}

% Slide 7: Significance, p-value, decision rule
\begin{frame}{Significance, $p$-value, and Decision Rule}
  \footnotesize
  \begin{itemize}\itemsep=0.12cm
    \item \textbf{Significance level ($\alpha$):} the maximum acceptable probability of a Type I Error, fixed \emph{before} observing the data (typically $0.10$, $0.05$, or $0.01$). \pause
    \item \textbf{$p$-value:} the probability of observing a result at least as extreme as the one obtained, assuming $H_0$ is true. \pause
    \item \textbf{Decision rule:} if $p\text{-value} \le \alpha$, reject $H_0$; otherwise, do not reject (which is \emph{not} the same as accepting it).
  \end{itemize}
\end{frame}

% Slide 8: Python Lab (Block 1)
\begin{frame}[fragile]{Python Lab: Type I Error Rate}
  \scriptsize
  Monte Carlo verification that the false-rejection rate equals $\alpha$ under a true $H_0$ (\texttt{07.01\_hypothesis\_testing\_basics.py}):
  {\lstset{basicstyle=\fontsize{4pt}{4.8pt}\selectfont\ttfamily,aboveskip=0.1em,belowskip=0.1em}
  \lstinputlisting[firstline=17, lastline=36]{../../code/07_pruebas_hipotesis/07.01_hypothesis_testing_basics.py}}
\end{frame}

% Slide 9: Python Lab (Block 2)
\begin{frame}[fragile]{Python Lab: Type II Error and Power Function}
  \scriptsize
  Analytic computation and simulation-based verification of $\beta$ and power:
  {\lstset{basicstyle=\fontsize{4pt}{4.8pt}\selectfont\ttfamily,aboveskip=0.1em,belowskip=0.1em}
  \lstinputlisting[firstline=39, lastline=62]{../../code/07_pruebas_hipotesis/07.01_hypothesis_testing_basics.py}}
\end{frame}

% Slide 10: Python Lab (Block 3)
\begin{frame}[fragile]{Python Lab: Sample Size for a Target Power}
  \scriptsize
  Computing the minimum $n$ and verifying the achieved power empirically:
  {\lstset{basicstyle=\fontsize{4pt}{4.8pt}\selectfont\ttfamily,aboveskip=0.1em,belowskip=0.1em}
  \lstinputlisting[firstline=70, lastline=92]{../../code/07_pruebas_hipotesis/07.01_hypothesis_testing_basics.py}}
\end{frame}

% Slide 11: Terminal Output
\begin{frame}[fragile]{Python Lab: Terminal Output}
  \scriptsize
  Standard output produced when running the Python lab script:
  \vspace{0.02cm}
  \begin{block}{Standard Console Output}
    \fontsize{4pt}{5pt}\selectfont
    \begin{verbatim}
=== Block 1: Type I Error Rate ===
Critical value z_crit (alpha=0.05): 1.6449
Empirical Type I error rate: 0.0503  (nominal alpha: 0.0500)

=== Block 2: Type II Error and Power ===
Analytic beta (mu_a=108): 0.1534, Analytic power: 0.8466
Empirical power via Monte Carlo: 0.8460

=== Block 3: Sample Size for Target Power ===
Exact n = 30.11 -> rounded up: n=31
Verification with n=31: empirical power = 0.9073 (target=0.90)
    \end{verbatim}
  \end{block}
\end{frame}


% Slide 16: Closing
\begin{frame}{Section Closing: Foundations of Hypothesis Testing}
  \begin{reflexion}
    Every hypothesis test is a decision procedure under uncertainty, explicitly calibrated by two opposing risks: $\alpha$ (rejecting too often) and $\beta$ (rejecting too rarely). Designing a good experiment means consciously choosing $\alpha$, $n$, and the desired power --- not leaving them to chance.
  \end{reflexion}

  \vspace{0.2cm}
  \pause
  \begin{block}{Key Takeaways}
    \begin{itemize}\itemsep=0.08cm
      \item \textbf{$H_0$ vs. $H_a$:} the known/claimed value goes in $H_0$; the suspicion or finding goes in $H_a$.
      \item \textbf{$\alpha$ and $\beta$:} complementary risks; reducing one (with $n$ fixed) generally increases the other.
      \item \textbf{Power and sample size:} $n=((Z_\alpha+Z_\beta)\sigma/(\mu_a-\mu_0))^2$ links experimental design to tolerated risk.
    \end{itemize}
  \end{block}
\end{frame}

% Slide 17: Modular Perspective
\begin{frame}{Modular Perspective --- The Next Challenge}
  \scriptsize
  \begin{retoclase}
    We now master the full logical framework of hypothesis testing. Section 07.02 applies this framework to $Z$ and $t$ tests for one and two means --- independent samples with equal or unequal variances (Welch), and paired samples --- extending the single-mean example to group comparisons.
  \end{retoclase}

  \vspace{0.08cm}
  \pause
  \begin{alertblock}{Recommended Preparation}
    \begin{itemize}\itemsep=0.06cm
      \item Review Student's $t$-distribution and its degrees of freedom (Section 05.04).
      \item Reflect on why pooling variances ($S_p^2$) requires the homoscedasticity assumption.
    \end{itemize}
  \end{alertblock}
\end{frame}

\end{document}
